\documentclass[11pt,a4paper]{article}
\usepackage[T1]{fontenc}
\usepackage[utf8]{inputenc}
\usepackage{lmodern}
\usepackage[french]{babel}
\usepackage[margin=2cm,headheight=14pt]{geometry}
\usepackage[most]{tcolorbox}
\usepackage{xcolor,tabularx,booktabs,array,enumitem,fancyhdr,amssymb,tikz}
\usetikzlibrary{arrows.meta,positioning}

\pagestyle{fancy}
\fancyhf{}
\lhead{\small Robotique -- Grade 6 -- Novembre 2026}
\rhead{\small Leçon 6 : Transformation et transmission de l'énergie}
\cfoot{\thepage}
\setlist{itemsep=1pt,topsep=2pt,leftmargin=*}
\renewcommand{\arraystretch}{1.35}
\newcolumntype{Y}{>{\raggedright\arraybackslash}X}
\newcommand{\cb}{$\square$\,}
\newcommand{\fl}{$\rightarrow$}
\newcommand{\lig}{\noindent\rule{\linewidth}{0.3pt}\\[5mm]}

\newtcolorbox{etape}[3]{enhanced,breakable,colback=#1!6,colframe=#1!70!black,
  coltitle=white,fonttitle=\bfseries,title={#2\hfill #3},boxrule=0.8pt,arc=2mm,
  before skip=8pt,after skip=8pt}
\newtcolorbox{astuce}[1][Différenciation et besoins particuliers]{breakable,
  colback=yellow!12,colframe=yellow!50!black,fonttitle=\bfseries\small,
  title={#1},boxrule=0.5pt,arc=1mm,before skip=6pt,after skip=6pt}
\newcommand{\entete}[2]{%
  \begin{tcolorbox}[colback=#1!10,colframe=#1!70!black,boxrule=0.8pt,arc=2mm]
  {\Large\bfseries #2}\hfill Nom : \rule{4cm}{0.3pt}\quad Date : \rule{2.2cm}{0.3pt}
  \end{tcolorbox}}

% Chaîne énergétique : #1 source, #2 transformation, #3 transmission, #4 utilisation
\newcommand{\chaine}[4]{%
\begin{center}
\begin{tikzpicture}[box/.style={draw,rounded corners,thick,minimum height=1.3cm,text width=3cm,align=center,font=\small},>={Stealth[length=3mm]},node distance=0.8cm]
\node[box,fill=orange!20] (a) {\textbf{Source}\\ #1};
\node[box,fill=blue!15,right=of a] (b) {\textbf{Transformation}\\ #2};
\node[box,fill=green!15,right=of b] (c) {\textbf{Transmission}\\ #3};
\node[box,fill=violet!15,right=of c] (d) {\textbf{Utilisation}\\ #4};
\draw[->,thick] (a)--(b); \draw[->,thick] (b)--(c); \draw[->,thick] (c)--(d);
\node[draw=red!70!black,dashed,rounded corners,font=\small,text=red!70!black,anchor=north east] (p) at ([yshift=-0.6cm]d.south east) {Pertes : une partie de l'énergie devient de la chaleur};
\draw[->,red!70!black,dashed] (d.south)--(d.south |- p.north);
\end{tikzpicture}
\end{center}}

\begin{document}

\begin{tcolorbox}[colback=teal!10,colframe=teal!70!black,arc=2mm,boxrule=1pt]
\centering
{\LARGE\bfseries Leçon 6 : Transformation et transmission de l'énergie}\\[2mm]
{\large Robotique -- Grade 6 -- Plan de cours de 2 heures}\\[1mm]
Enseignant : Francis Ki \quad|\quad Semaine du 23 au 27 novembre 2026
\end{tcolorbox}

%==================================================================
\section{Informations générales}
\begin{tabularx}{\linewidth}{>{\bfseries}p{3.6cm}Y}
\toprule
Titre & Leçon 6 : Transformation et transmission de l'énergie\\
Classe / effectif & Grade 6, 11 à 20 élèves, niveaux hétérogènes, élèves à besoins particuliers\\
Durée & 2 heures (120 min)\\
Lieu & Salle informatique, un poste par élève\\
Matériel et logiciels &
\begin{itemize}
 \item Un poste par élève avec navigateur web à jour et accès internet.
 \item Vidéoprojecteur ou tableau interactif pour l'enseignant.
 \item Simulation PhET «~Formes et changements d'énergie~» (phet.colorado.edu, gratuite, en français), déjà utilisée à la leçon 5. \textbf{À tester sur un poste élève avant la séance.}
 \item Diapositives (Google Slides ou LibreOffice Impress) pour dessiner les chaînes énergétiques, avec des formes et des flèches.
 \item Fiches numériques partagées en lecture seule : chaque élève crée sa copie et la renomme avec son prénom.
 \item Facultatif : une lampe de poche démontable et un petit moteur avec une pile pour la démonstration.
\end{itemize}\\
Pré-requis &
Leçon 5 : connaître les formes d'énergie (électrique, chimique, lumineuse, thermique, mécanique, sonore) et les quatre effets (lumière, chaleur, mouvement, son).\\
\bottomrule
\end{tabularx}

\begin{astuce}[Préparation avant le cours (environ 35 min)]
\begin{enumerate}
 \item Tester la simulation PhET sur un poste élève. Prévoir un plan B sans simulation (schémas de l'annexe A).
 \item Préparer les fichiers numériques des annexes B à E (une copie par élève) avec des cases vides pour les chaînes.
 \item Créer le formulaire de sortie (annexe G) dans Google Forms ou imprimer l'annexe.
 \item Reconstituer les groupes A, B, C à partir de l'exit ticket de la leçon 5.
\end{enumerate}
\end{astuce}

%==================================================================
\section{Objectifs d'apprentissage}
\begin{tcolorbox}[colback=gray!8,colframe=gray!60!black,arc=1mm,boxrule=0.5pt]
\textbf{Objectifs du programme}
\begin{itemize}
 \item Comprendre comment l'énergie se transforme dans un système technique.
 \item Décrire la chaîne énergétique d'un objet simple.
\end{itemize}
\textbf{À la fin de la séance, l'élève est capable\,:}
\begin{itemize}
 \item d'expliquer qu'une énergie change de forme dans un objet (par exemple électrique en lumineuse)\,;
 \item de distinguer \emph{transformer} (changer de forme) et \emph{transmettre} (faire voyager)\,;
 \item de compléter, puis de construire, la chaîne énergétique d'un objet simple : source, transformation, transmission, utilisation.
\end{itemize}
\textbf{Compétences numériques\,:} dessiner un schéma avec formes et flèches, manipuler une simulation, rendre son travail.
\end{tcolorbox}

%==================================================================
\section{Déroulé minuté (120 min)}

\begin{tabularx}{\linewidth}{>{\bfseries}p{2.2cm}p{2.6cm}Y p{2.8cm}}
\toprule
Minutes & Étape & Contenu & Organisation\\
\midrule
0 -- 10 & Accroche & La lampe de poche démontée & Classe entière, sans écran\\
10 -- 30 & Mini-leçon & Transformer, transmettre et la chaîne énergétique & Classe entière\\
30 -- 60 & Pratique guidée & Un système dans la simulation, puis deux objets en binômes & Collectif puis binômes\\
60 -- 95 & Pratique autonome & Trois niveaux de besoin & Individuel\\
95 -- 110 & Mise en commun & Valoriser chaque groupe & Classe entière\\
110 -- 120 & Exit ticket & Trois questions sur poste & Individuel\\
\bottomrule
\end{tabularx}

%------------------------------------------------------------------
\begin{etape}{orange}{1. Accroche : la lampe de poche démontée}{0 -- 10 min}
\textbf{But\,:} faire penser au trajet de l'énergie, de la pile à la lumière.\\
\textbf{Écrans éteints.}
\begin{enumerate}
 \item (4 min) Montrer une lampe de poche (ouverte si possible : pile, interrupteur, ampoule). Question : \emph{Que se passe-t-il dans cet objet quand on appuie ? Par où passe l'énergie ?}
 \item (3 min) Poser l'ampoule allumée sur la main un instant (sans danger si elle est de faible puissance, sinon utiliser un radiateur ou un ordinateur) : \emph{Pourquoi l'ampoule chauffe-t-elle alors qu'on voulait de la lumière ?}
 \item (3 min) Question-problème : \textbf{«~Que devient l'énergie de la pile quand la lampe est allumée ?~»}
\end{enumerate}
\textit{Sans lampe réelle, projeter une photo légendée ou un schéma de lampe de poche.}
\end{etape}

%------------------------------------------------------------------
\begin{etape}{blue}{2. Mini-leçon : transformer, transmettre, chaîne énergétique}{10 -- 30 min}
\textbf{Support\,:} schéma construit pas à pas au tableau ou projeté (annexe A). Quatre temps de 5 min.

\medskip
\textbf{Temps 1 : l'énergie change de forme.} L'énergie ne disparaît pas : elle se transforme. Exemples : une pile transforme de l'énergie chimique en électrique ; une lampe transforme de l'électrique en lumineuse ; un panneau solaire transforme de la lumineuse en électrique.

\textbf{Temps 2 : l'énergie voyage.} \emph{Transmettre}, c'est faire voyager l'énergie d'un endroit à un autre sans changer sa forme.

\begin{center}
\begin{tabularx}{\linewidth}{>{\bfseries}p{3.2cm}Y Y}
\toprule
 & Transformer & Transmettre\\
\midrule
Définition & L'énergie change de forme & L'énergie va d'un point à un autre\\
Exemple & Ampoule : électrique \fl{} lumineuse & Fils : l'énergie électrique va de la pile à l'ampoule\\
Autre exemple & Moteur : électrique \fl{} mécanique & Chaîne de vélo : le mouvement va des pédales à la roue\\
\bottomrule
\end{tabularx}
\end{center}

\textbf{Temps 3 : la chaîne énergétique.} Pour décrire un objet, on suit l'énergie de sa source jusqu'à son utilisation.

\chaine{Pile (chimique)}{Pile : chimique \fl{} électrique}{Fils (électrique)}{Ampoule : électrique \fl{} lumière}
\textit{Exemple : la lampe de poche. Pour un autre objet, certaines cases peuvent être vides (un fil peut manquer, par exemple) ou la transformation peut avoir lieu à la fin.}

\textbf{Temps 4 : les pertes.} Dans chaque transformation, une partie de l'énergie devient de la chaleur qu'on n'utilise pas. C'est pourquoi l'ampoule chauffe. Retour à la question-problème : \textit{l'énergie de la pile devient de la lumière utile et de la chaleur perdue.}

\begin{astuce}
Dessiner la chaîne avec les mêmes couleurs de cases pour la source, la transformation, la transmission et l'utilisation, sur tous les supports. Lire les schémas à voix haute en suivant les flèches avec le doigt.
\end{astuce}
\end{etape}

%------------------------------------------------------------------
\begin{etape}{green}{3. Pratique guidée}{30 -- 60 min}
\textbf{Fiche\,:} annexe B (version numérique : fichier «~L6\_Guidee~»).

\medskip
\textbf{3a. Exemple collectif sur poste : un système dans la simulation (12 min).} Tout le monde ouvre la simulation PhET «~Formes et changements d'énergie~».
\begin{enumerate}
 \item Construire un système : une source (par exemple le robinet ou le cycliste), un convertisseur (roue ou générateur) et un utilisateur (ampoule ou ventilateur). Afficher les symboles d'énergie.
 \item Suivre les symboles de la source à l'utilisateur et compléter ensemble le schéma B1 : forme d'énergie à chaque étape, transformation ou transmission, effet utile, perte.
\end{enumerate}

\textbf{3b. Deux objets en binômes (18 min).} Objet 1 : un grille-pain. Objet 2 : un ventilateur. Remplir le schéma B2 pour chacun. Alternance des rôles à mi-temps.

\textbf{Étayage progressif de l'aide.}
\begin{center}
\begin{tabularx}{\linewidth}{>{\bfseries}p{2.4cm}Y p{4.6cm}}
\toprule
Niveau d'aide & Contenu & Qui intervient\\
\midrule
Aide 1 & Chaîne type de la lampe (annexe A) et tableau des formes d'énergie de la leçon 5 & Le binôme seul\\
Aide 2 & Questions guides : \emph{D'où vient l'énergie ? Où change-t-elle de forme ? Qu'est-ce qui est perdu ?} & Le binôme avec la carte\\
Aide 3 & L'enseignant remplit une case avec le binôme, qui termine & L'enseignant\\
\bottomrule
\end{tabularx}
\end{center}

\begin{astuce}
Si la simulation est trop rapide ou trop riche pour un élève, utiliser seulement le schéma de la lampe. Proposer un schéma pré-rempli à compléter au lieu d'un schéma vierge. Un binôme tuteur peut aider à manipuler la simulation.
\end{astuce}
\end{etape}

%------------------------------------------------------------------
\begin{etape}{violet}{4. Pratique autonome différenciée}{60 -- 95 min}
\textbf{Consigne commune\,:} chaque élève ouvre le fichier de son groupe, le renomme avec son prénom et travaille seul. L'enseignant passe d'abord auprès du groupe A pendant 10 min.

\medskip
\begin{tabularx}{\linewidth}{>{\bfseries}p{3cm}Y p{3.2cm}}
\toprule
Groupe & Tâche sur poste & Annexe\\
\midrule
A : en dessous du niveau &
Trois chaînes à compléter avec une banque de mots (lampe de poche, ventilateur, panneau solaire) ; relier les mots «~transformer~» et «~transmettre~» à leur définition. & Annexe C\\
B : au niveau &
Décrire la chaîne énergétique de trois objets (bouilloire, vélo, voiture jouet à pile) avec un schéma sur diapositive : formes d'énergie, transformation, transmission, effet utile, perte. & Annexe D\\
C : au-dessus du niveau &
Mission «~Réduire les pertes~» : comparer deux ampoules, imaginer une voiture solaire jouet et dessiner sa chaîne énergétique, puis relier avec un robot (batterie, moteur, roues). & Annexe E\\
\bottomrule
\end{tabularx}

\medskip
\textbf{Plan B si un élève termine tôt\,:} retourner dans la simulation et enchaîner deux convertisseurs, puis décrire ce qu'on observe.
\end{etape}

%------------------------------------------------------------------
\begin{etape}{teal}{5. Mise en commun}{95 -- 110 min}
\begin{enumerate}
 \item (2 min) Retour au calme, écrans en veille.
 \item (9 min) Trois passages courts de 3 min :
 \begin{itemize}
  \item \textbf{Groupe A} : un élève lit à voix haute une chaîne complétée. Valoriser la lecture dans l'ordre des flèches.
  \item \textbf{Groupe B} : un binôme projette un schéma et explique où l'énergie est transformée. Valoriser la précision.
  \item \textbf{Groupe C} : présentation de la voiture solaire et du lien avec le robot. Valoriser la créativité.
 \end{itemize}
 \item (4 min) Revenir à la question-problème et conclure : \textit{l'énergie ne disparaît pas, elle change de forme et une partie est perdue sous forme de chaleur.} Annoncer que la chaîne énergétique reviendra pour le projet robotique.
\end{enumerate}
\end{etape}

\begin{etape}{gray}{6. Clôture : exit ticket}{110 -- 120 min}
Formulaire en ligne ou papier (annexe G) :
\begin{enumerate}
 \item Dans une lampe de poche, quelles formes d'énergie rencontre-t-on ?
 \item Vrai ou faux : l'énergie disparaît quand on éteint la lampe. Explique.
 \item Auto-évaluation : \emph{Je maîtrise / Je progresse / J'ai besoin d'aide.}
\end{enumerate}
\end{etape}

%==================================================================
% ANNEXES
%==================================================================
\newpage
\section*{Annexe A : Support résumé de la mini-leçon}
\entete{blue}{Transformation et transmission de l'énergie}

\begin{tcolorbox}[colback=yellow!10,colframe=yellow!50!black]
\textbf{L'énergie ne disparaît pas : elle change de forme (transformation) et elle voyage (transmission).}
\end{tcolorbox}

\begin{tabularx}{\linewidth}{>{\bfseries}p{3.2cm}Y Y}
\toprule
 & Transformer & Transmettre\\
\midrule
Définition & L'énergie change de forme & L'énergie va d'un point à un autre\\
Exemple & Ampoule : électrique \fl{} lumineuse & Fils, chaîne de vélo\\
\bottomrule
\end{tabularx}

\chaine{Pile (chimique)}{Pile : chimique \fl{} électrique}{Fils (électrique)}{Ampoule : électrique \fl{} lumière}

\begin{tcolorbox}[colback=green!8,colframe=green!50!black]
\textbf{Phrase modèle :} La \underline{\hspace{2cm}} reçoit de l'énergie \underline{\hspace{2cm}}, la transforme en \underline{\hspace{2cm}} et en \underline{\hspace{2cm}} (perdue).
\end{tcolorbox}

%==================================================================
\newpage
\section*{Annexe B : Fiche de pratique guidée}
\entete{green}{Pratique guidée : la chaîne énergétique}

\subsection*{B1. Exemple collectif dans la simulation}
Système construit : source \dotfill{} \quad convertisseur \dotfill{} \quad utilisateur \dotfill

\begin{tabularx}{\linewidth}{>{\bfseries}p{3.4cm}Y Y Y Y}
\toprule
Étape & Source & Transformation & Transmission & Utilisation\\
\midrule
Forme d'énergie & & & & \\
Effet utile & & & & \\
Perte & & & & \\
\bottomrule
\end{tabularx}

\subsection*{B2. En binôme : deux objets}
\begin{tabularx}{\linewidth}{>{\bfseries}p{3cm}Y Y Y Y}
\toprule
Objet & Source & Transformation & Transmission & Utilisation et perte\\
\midrule
Grille-pain & & & & \\
Ventilateur & & & & \\
\bottomrule
\end{tabularx}

\begin{tcolorbox}[colback=yellow!10,colframe=yellow!50!black,title={Cartes d'aide à utiliser dans l'ordre},fonttitle=\bfseries\small]
\textbf{Aide 1 :} regarde la chaîne de la lampe de poche (annexe A).\\
\textbf{Aide 2 :} D'où vient l'énergie ? Où change-t-elle de forme ? Qu'est-ce qui est perdu ?\\
\textbf{Aide 3 :} demande à l'enseignant de remplir une case avec toi.
\end{tcolorbox}

%==================================================================
\newpage
\section*{Annexe C : Niveau A (en dessous du niveau)}
\entete{orange}{Niveau A : je complète les chaînes}

\textbf{Banque de mots :} pile, fils, ampoule, électrique, chimique, lumière, chaleur, soleil, panneau solaire, moteur, hélice, mouvement.

\subsection*{Activité 1 : complète avec la banque de mots (glisse les étiquettes)}
\begin{tabularx}{\linewidth}{>{\bfseries}p{3.2cm}Y Y Y Y}
\toprule
Objet & Source & Transformation & Transmission & Utilisation\\
\midrule
Lampe de poche & \dotfill & \dotfill{} & fils & lumière\\
Ventilateur & prise & moteur & \dotfill & \dotfill\\
Panneau solaire & \dotfill & \dotfill & fils & électricité\\
\bottomrule
\end{tabularx}

\subsection*{Activité 2 : relie}
\begin{tabularx}{\linewidth}{Y Y}
\toprule
\textbf{Mot} & \textbf{Définition}\\
\midrule
Transformer & L'énergie voyage d'un point à un autre\\
Transmettre & L'énergie change de forme\\
\bottomrule
\end{tabularx}
Réponse : Transformer \fl{} \dotfill\quad Transmettre \fl{} \dotfill

\subsection*{Activité 3 : complète}
La lampe transforme l'énergie \dotfill{} en \dotfill{} et en un peu de \dotfill.

%==================================================================
\newpage
\section*{Annexe D : Niveau B (au niveau)}
\entete{blue}{Niveau B : je décris des chaînes}

Pour chaque objet, dessine sur une diapositive la chaîne énergétique (quatre cases et des flèches) puis complète le tableau.

\begin{tabularx}{\linewidth}{>{\bfseries}p{3.2cm}Y Y Y Y}
\toprule
Objet & Source (forme) & Transformation & Transmission & Effet utile et perte\\
\midrule
Bouilloire électrique & & & & \\
Vélo & & & & \\
Voiture jouet à pile & & & & \\
\bottomrule
\end{tabularx}

\subsection*{Question de réflexion}
Pourquoi un ordinateur chauffe-t-il alors qu'on veut seulement de l'image et du son ?\\
\lig \lig

%==================================================================
\newpage
\section*{Annexe E : Niveau C (au-dessus du niveau)}
\entete{violet}{Niveau C : réduire les pertes}

\subsection*{Mission 1 : comparer deux ampoules}
Une ampoule à filament chauffe beaucoup. Une ampoule LED chauffe très peu pour la même lumière. Dessine les deux chaînes sur une diapositive. Laquelle gaspille le moins d'énergie ? Explique.\\
\lig \lig

\subsection*{Mission 2 : conçois une voiture solaire jouet}
Dessine sa chaîne énergétique : source, transformation, transmission, utilisation, et indique où se trouvent les pertes.\\
Avec quels matériaux (leçon 4) construirais-tu la caisse ?\\
\lig

\subsection*{Mission 3 : relie avec un robot}
Un robot a une batterie, un moteur et des roues. Écris sa chaîne énergétique.\\
Source : \dotfill{} \quad Transformation : \dotfill{} \quad Transmission : \dotfill{} \quad Utilisation : \dotfill

\textbf{Défi bonus :} dans la simulation, enchaîne deux convertisseurs et décris ce qui change.

%==================================================================
\newpage
\section*{Annexe F : Corrections et attendus (enseignant)}

\subsection*{Annexe B : pratique guidée}
\begin{tabularx}{\linewidth}{>{\bfseries}p{3cm}Y Y Y Y}
\toprule
Objet & Source & Transformation & Transmission & Utilisation et perte\\
\midrule
Grille-pain & Prise (électrique) & Résistance : électrique \fl{} thermique & Fils (électrique) & Chaleur pour griller ; un peu de lumière perdue\\
Ventilateur & Prise (électrique) & Moteur : électrique \fl{} mécanique & Axe et pales (mécanique) & Mouvement de l'air ; chaleur et bruit perdus\\
\bottomrule
\end{tabularx}
Simulation (exemple) : robinet \fl{} roue \fl{} générateur \fl{} ampoule, avec mécanique \fl{} électrique \fl{} lumineuse et thermique.

\subsection*{Annexe C : niveau A}
Lampe : pile, ampoule. Ventilateur : fils, mouvement. Panneau : soleil, panneau solaire. Activité 2 : Transformer \fl{} l'énergie change de forme ; Transmettre \fl{} l'énergie voyage. Activité 3 : chimique, lumineuse, chaleur.

\subsection*{Annexe D : niveau B}
\begin{tabularx}{\linewidth}{>{\bfseries}p{3.2cm}Y Y Y Y}
\toprule
Objet & Source & Transformation & Transmission & Effet utile / perte\\
\midrule
Bouilloire & Prise & Électrique \fl{} thermique & Fils & Eau chaude / chaleur de l'air\\
Vélo & Cycliste (chimique) & Muscles : chimique \fl{} mécanique & Pédales, chaîne & Mouvement / chaleur des frottements\\
Voiture à pile & Pile & Moteur : électrique \fl{} mécanique & Fils, engrenages & Mouvement / chaleur, bruit\\
\bottomrule
\end{tabularx}
Réflexion : l'électronique transforme de l'énergie électrique en lumière, en son et en chaleur ; la chaleur est une perte qui chauffe l'ordinateur.

\subsection*{Annexe E : niveau C}
Ampoule à filament : beaucoup de chaleur perdue ; LED : moins de pertes, donc plus économe. Voiture solaire : soleil \fl{} panneau \fl{} fils \fl{} moteur \fl{} roues. Robot : batterie (chimique), moteur (électrique \fl{} mécanique), transmission par fils et axes, roues.\\
\textbf{Critères de réussite :} deux chaînes correctes, une comparaison argumentée, une chaîne de robot complète.

%==================================================================
\newpage
\section*{Annexe G : Exit ticket et grille d'observation}
\entete{gray}{Exit ticket : transformation et transmission}

\begin{enumerate}
 \item Dans une lampe de poche, quelles formes d'énergie rencontre-t-on ?\\ \lig
 \item Vrai ou faux : l'énergie disparaît quand on éteint la lampe. Explique.\\ \lig
 \item Je me sens : \cb Je maîtrise \quad \cb Je progresse \quad \cb J'ai besoin d'aide
\end{enumerate}

\textbf{Attendus :} 1. Chimique (pile), électrique (fils), lumineuse (lumière) et thermique (chaleur). 2. Faux : l'énergie n'est plus transformée en lumière, mais elle n'a pas disparu (elle reste dans la pile, ou devient chaleur).

\bigskip
\subsection*{Grille d'observation (enseignant)}
\begin{tabularx}{\linewidth}{Y p{2.4cm} p{2.4cm} p{2.4cm} Y}
\toprule
Élève & Transformation & Transmission & Chaîne complète & Groupe suivant\\
\midrule
 & & & & \\ & & & & \\ & & & & \\ & & & & \\ & & & & \\
\bottomrule
\end{tabularx}

\end{document}
